\documentclass[aspectratio=169]{beamer}

\usepackage{amsmath,amssymb}
\usepackage{booktabs}
\usepackage{tabularx}
\usepackage{tikz}
\usetikzlibrary{arrows.meta,backgrounds,calc,positioning,shapes.multipart}
\usepackage[main=english,polish]{babel}
\usepackage{ragged2e}
\usepackage{csquotes}
\usepackage[
  backend=biber,
  style=authoryear,
  maxcitenames=2,
  doi=false,
  isbn=false,
  url=true
]{biblatex}
\addbibresource{refs.bib}

\definecolor{cookieDemoBlue}{HTML}{356AE6}
\usetheme[
  accent=cookieDemoBlue,
  progressbar=foot,
  sectionpage=progressbar,
  subsectionpage=progressbar,
  numbering=fraction,
  numberrail=subsection,
  block=fill,
  notes=hide,
  toc=aftertitle,
  closing=contact
]{cookie}

\lstset{style=cookie}

\title[Cookie]{Cookie Beamer Theme}
\subtitle{A LuaLaTeX Beamer theme with modern visuals}
\author[Mars]{SeniorMars}
\cookieemail{seniormars@rice.edu}
\institute{Mathematics}
\cookieuni{Rice University}
\cookielocation{Houston, Texas}
\date{June 21, 2026}
\cookietagline{A small theme for modern slides}
\cookietitleimage{assets/cookie-background.png}
\cookieclosingtitle{Thank you}
\cookieclosingtext{Questions, forks, and bug reports are welcome.}
\cookieclosingqr{https://example.org/cookie-beamer}
\cookieclosingqrlabel{Replace with your project URL}

\titlegraphic{%
  \begin{tikzpicture}[scale=.70]
    \fill[white,opacity=.94] (0,0) circle[radius=1.15];
    \fill[cookieDemoBlue!82!black] (-.43,.40) circle[radius=.12];
    \fill[cookieDemoBlue!82!black] (.36,.50) circle[radius=.10];
    \fill[cookieDemoBlue!82!black] (.52,-.22) circle[radius=.13];
    \fill[cookieDemoBlue!82!black] (-.38,-.43) circle[radius=.09];
    \fill[cookieDemoBlue!82!black] (.02,.02) circle[radius=.11];
  \end{tikzpicture}%
}

\newsavebox{\cookieDemoPenguin}
\savebox{\cookieDemoPenguin}{\cookiepenguin[scale=.31]}
\newcommand{\cookiedemomathcomparisonbody}{%
  \footnotesize
  \setlength{\abovedisplayskip}{.24em}%
  \setlength{\belowdisplayskip}{.24em}%
  \begin{columns}[T,onlytextwidth]
    \column{.18\textwidth}
      \centering
      \cookiekicker{Lining}\par
      \vspace{.45em}
      1234567890

    \column{.18\textwidth}
      \centering
      \cookiekicker{Oldstyle}\par
      \vspace{.45em}
      \oldstylenums{1234567890}

    \column{.22\textwidth}
      \centering
      \cookiekicker{Accents}\par
      \vspace{.45em}
      $\hat{x},\;\check{x},\;\tilde{a},\;\bar{a},\;\dot{y},\;\ddot{y}$

    \column{.36\textwidth}
      \centering
      \cookiekicker{Differentials}\par
      \vspace{.45em}
      $\displaystyle \iiint f(x,y,z)\,\mathrm dx\,\mathrm dy\,\mathrm dz$
  \end{columns}

  \vspace{.30em}
  \begin{columns}[c,onlytextwidth]
    \column{.47\textwidth}
      \[
        \frac{1}{\displaystyle 1+
          \frac{1}{\displaystyle 2+
          \frac{1}{\displaystyle 3+x}}}
        +
        \frac{1}{1+\frac{1}{2+\frac{1}{3+x}}}
      \]

    \column{.47\textwidth}
      \[
        F:\left|\begin{array}{ccc}
          F''_{xx} & F''_{xy} & F'_x \\
          F''_{yx} & F''_{yy} & F'_y \\
          F'_x     & F'_y     & 0
        \end{array}\right|=0
      \]
  \end{columns}

  \vspace{.02em}
  \begin{columns}[c,onlytextwidth]
    \column{.30\textwidth}
      \[
        \mathop{\int\!\!\!\int}_{\mathbf{x}\in\mathbb R^2}
        \!\langle \mathbf{x},\mathbf{y}\rangle\,\mathrm d\mathbf{x}
      \]

    \column{.34\textwidth}
      \[
        \overline{\overline{a\alpha}^2+\underline{b\beta}
        +\overline{\overline{d\delta}}}
      \]

    \column{.30\textwidth}
      \[
        \left]0,1\right[+\lceil x\rfloor-\langle x,y\rangle
      \]
  \end{columns}

  \vspace{.05em}
  \begin{columns}[c,onlytextwidth]
    \column{.44\textwidth}
      \[
        \begin{aligned}
          e^x &\approx 1+x+x^2/2! \\
              &\quad{}+x^3/3!+x^4/4!
        \end{aligned}
      \]

    \column{.50\textwidth}
      \[
        \binom{n+1}{k}=\binom{n}{k}+\binom{n}{k-1}
      \]
  \end{columns}

  \vspace{-.15em}
  {\scriptsize
  \[
    \begin{aligned}
      p(R,\phi)
      &\sim \int_{-\infty}^{\infty}
      \frac{
        \tilde{W}_n(\gamma)
        \exp\!\left[
          \imath R/a
          \left(\sqrt{k^2a^2-\gamma^2}\cos\phi\right)
        \right]
      }{
        \left(k^2a^2-\gamma^2\right)^{3/4}
      }\\[-.15em]
      &\qquad{}\times
      \frac{1}{
        {H_n^{\prime}}_n^{(1)}
        \left(\sqrt{k^2a^2-\gamma^2}\right)
      }\,d\gamma .
    \end{aligned}
  \]
  }%
}
\begin{document}

\maketitle

\section{Design system}

\begin{frame}{A theme for modern talks}
  \begin{columns}[T,onlytextwidth]
    \column{.31\textwidth}
      \cookiekicker{Awesome-style frame}
      \vspace{.55em}

      Number rails, angled title art, and section slides are built in.

    \column{.31\textwidth}
      \cookiekicker{Metropolis defaults}
      \vspace{.55em}

      Progress bars, block styles, and fonts are set from theme options.

    \column{.31\textwidth}
      \cookiekicker{LuaLaTeX setup}
      \vspace{.55em}

      Uses installed system fonts, Unicode math, and no private package paths.
  \end{columns}

  \vfill
  \cookiebadge{single .sty} \hspace{.35em}
  \cookiebadge{LuaLaTeX} \hspace{.35em}
  \cookiebadge{old-deck friendly}
\end{frame}

\begin{frame}{Navigation stays visible}{Long titles and subtitles fit the header}
  \begin{columns}[T,onlytextwidth]
    \column{.58\textwidth}
      \begin{itemize}
        \item The rail shows the current section, or section.subsection when needed.
        \item The footer keeps author, title, section, and frame count out of the slide body.
        \item Section pages mark turns in the talk without adding clutter to every frame.
      \end{itemize}

    \column{.36\textwidth}
      \begin{tikzpicture}[node distance=5mm]
        \node[cookie accent node,minimum width=3.6cm] (a) {Section};
        \node[cookie node,below=of a,minimum width=3.6cm] (b) {Subsection};
        \node[cookie node,below=of b,minimum width=3.6cm] (c) {Frame};
        \draw[cookie arrow] (a) -- (b);
        \draw[cookie arrow] (b) -- (c);
      \end{tikzpicture}
  \end{columns}
  \note{
    This frame intentionally uses a long title and a subtitle. The title area
    should grow without colliding with the divider.
  }
\end{frame}

\begin{frame}[label=lists]{Lists and local text}{Nested bullets, numbered steps, and Polish characters}
  \begin{columns}[T,onlytextwidth]
    \column{.48\textwidth}
      \begin{itemize}
        \item State the main point first.
        \begin{itemize}
          \item Add the detail that changes the argument.
          \begin{itemize}
            \item Keep the third level short.
          \end{itemize}
        \end{itemize}
        \item Return to the main thread.
      \end{itemize}

    \column{.43\textwidth}
      \begin{enumerate}
        \item Define the object.
        \item Prove the useful property.
        \item Use it in the example.
      \end{enumerate}
  \end{columns}

  \vspace{.75em}
  \justifying\small
  {\foreignlanguage{polish}{Tutaj jest krótki tekst po polsku z polskimi znakami: zażółć gęślą jaźń.}}

  \vspace{.35em}
  The quick, brown fox jumps over a lazy dog.
\end{frame}

\subsection{Components}

\begin{frame}[fragile]{Blocks and emphasis}
  \begin{columns}[T,onlytextwidth]
    \column{.48\textwidth}
      \begin{block}{A plain block}
        Use this for assumptions, notation, or a short claim.
      \end{block}

      \begin{exampleblock}{An example block}
        Use this for computations, test cases, or constructions.
      \end{exampleblock}

    \column{.48\textwidth}
      \begin{alertblock}{An alert block}
        Use this for a warning, failed case, or decision point.
      \end{alertblock}

      \begin{block}{Standard syntax}
        Cookie uses Beamer's normal \verb|\begin{block}{Title}| form.
      \end{block}
  \end{columns}
\end{frame}

\begin{frame}{Definitions, theorems, and proofs}
  \begin{definition}[Divisibility]
    For integers $a,b$, write $a\mid b$ exactly when
    $\exists c\in\mathbb Z$ such that $ac=b$.
  \end{definition}

  \begin{theorem}
    Every integer divides zero: $\forall a\in\mathbb Z,\; a\mid 0$.
  \end{theorem}

  \begin{proof}
    Choose $c=0$. Then $a\cdot c=a\cdot0=0$ for every integer $a$.
  \end{proof}
\end{frame}

\begin{frame}{Mathematics and numerals}
  \cookiedemomathcomparisonbody
\end{frame}

\begin{frame}[fragile]{Code is part of the slide}
  \begin{columns}[T,onlytextwidth]
    \column{.53\textwidth}
      \begin{cookiecode}[title=Scalar C]{C}
void square4(float *x) {
  for (int i = 0; i < 4; ++i) {
    x[i] = x[i] * x[i];
  }
}
      \end{cookiecode}

    \column{.43\textwidth}
      \begin{block}{What the theme sets}
        \begin{itemize}
          \item Noto Sans Mono when it is installed
          \item Line wrapping and padding inside the box
          \item Colors for keywords, strings, and comments
        \end{itemize}
      \end{block}

      \cookiebadge{listings} \hspace{.25em}\cookiebadge{no shell escape}
  \end{columns}
\end{frame}

\begin{frame}[fragile]{Longer code listing}{The title wraps without changing the code box}
  \begin{cookiecode}[title=Sleep sort sketch]{C}
#include <stdio.h>
#include <stdlib.h>
#include <unistd.h>

int main(int argc, char **argv) {
  while (--argc > 1 && !fork())
    sleep(atoi(argv[argc]));

  puts(argv[argc]);
  return 0;
}
  \end{cookiecode}
\end{frame}

\begin{frame}{Projected gradient descent}{Line numbers and comments stay readable}
  \begin{algorithm}[H]
    \caption{Projected gradient descent}
    \KwIn{Objective $f$, feasible set $C$, initial point $x_0$, step sizes $\eta_t$}
    \KwOut{Approximate minimizer $x_t$}
    \For{$t = 0,1,\ldots,T-1$}{
      $g_t \leftarrow \nabla f(x_t)$\tcp*{gradient at the current iterate}
      $y_{t+1} \leftarrow x_t - \eta_t g_t$\;
      $x_{t+1} \leftarrow \Pi_C(y_{t+1})$\tcp*{project back onto $C$}
      \If{$\lVert x_{t+1}-x_t\rVert < \varepsilon$}{
        \Return $x_{t+1}$\;
      }
    }
    \Return $x_T$\;
  \end{algorithm}
\end{frame}

\begin{frame}{RealBoost sketch}{A compact algorithm with math updates}
  \begin{algorithm}[H]
    \caption{RealBoost training loop}
    \KwIn{Training data $(x_i,y_i)_{i=1}^m$, rounds $T$}
    \KwOut{Classifier $F$}
    Initialize weights $w_i \gets 1/m$\;
    \For{$t=1,\ldots,T$}{
      Fit $f_t(x)\in\mathbb R$ using the current weights\;
      Set $Z_t \gets \sum_{i=1}^m w_i\exp(-f_t(x_i)y_i)$\;
      Update $w_i \gets w_i\exp(-f_t(x_i)y_i)/Z_t$\tcp*{normalize}
    }
    \Return $F(x)=\sum_{t=1}^{T} f_t(x)$\;
  \end{algorithm}
\end{frame}

\begin{frame}{Two trajectories side by side}
  \begin{columns}[T,onlytextwidth]
    \column{.29\textwidth}
      \begin{cookiecard}[title=Minimum fuel]
        \centering
        \begin{tikzpicture}[x=.47cm,y=.032cm]
          \draw[->,cookieMuted] (0,0) -- (6.3,0);
          \draw[->,cookieMuted] (0,0) -- (0,92);
          \draw[cookieAccent,line width=1.15pt,smooth]
            plot coordinates {(0.2,82)(1,64)(2,42)(3,27)(4.2,18)(5.3,12)(6,10)};
          \foreach \x/\y in {.7/72,1.5/53,2.4/35,3.4/23,4.5/16,5.5/11}
            \draw[-{Latex[length=1.3mm]},cookieInk!72] (\x,\y+7) -- (\x+.32,\y);
        \end{tikzpicture}

        \vspace{.35em}
        \footnotesize A smoother curve uses less fuel.
      \end{cookiecard}

    \column{.36\textwidth}
      \vspace{1.05em}
      \cookiekicker{Same dynamics}
      \vspace{.65em}

      \textbf{Left:} minimize fuel.

      \vspace{.85em}
      \textbf{Right:} minimize time.

      \vspace{.85em}
      The middle column names the tradeoff; the figures do the rest.

    \column{.29\textwidth}
      \begin{cookiecard}[title=Minimum time]
        \centering
        \begin{tikzpicture}[x=.47cm,y=.032cm]
          \draw[->,cookieMuted] (0,0) -- (6.3,0);
          \draw[->,cookieMuted] (0,0) -- (0,92);
          \draw[cookieAlert,line width=1.15pt,smooth]
            plot coordinates {(0.2,84)(.8,73)(1.5,56)(2.1,38)(2.8,23)(3.7,14)(5,10)(6,9)};
          \foreach \x/\y in {.55/77,1.2/64,1.8/46,2.5/29,3.3/18,4.3/12}
            \draw[-{Latex[length=1.3mm]},cookieInk!72] (\x,\y+7) -- (\x+.32,\y);
        \end{tikzpicture}

        \vspace{.35em}
        \footnotesize A steeper curve reaches the target sooner.
      \end{cookiecard}
  \end{columns}
\end{frame}

{
\setbeamercolor{normal text}{fg=white,bg=}
\setbeamercolor{frametitle}{fg=white,bg=}
\setbeamercolor{framesubtitle}{fg=white!78,bg=}
\setbeamercolor{footline}{fg=white!78,bg=cookieInk!72}
\begin{frame}[
  bgimage=assets/cookie-background.png,
  bgfit=cover,
  bgoverlay=cookieInk,
  bgoverlayopacity=.56
]{Frame-local background images}{Full canvas with readable text}
  \color{white}
  \begin{columns}[T,onlytextwidth]
    \column{.57\textwidth}
      \Large
      Use a photo, diagram, texture, or generated image on one frame without setting a deck-wide background.

      \vspace{.8em}
      \normalsize
      Set the wash color and opacity for contrast. Use local colors when the frame needs light text.

    \column{.38\textwidth}
      \begin{cookiecard}[title=Frame syntax]
        \ttfamily\small
        \textbackslash begin\{frame\}[\par
        \quad bgimage=photo.jpg,\par
        \quad bgoverlay=cookieInk,\par
        \quad bgoverlayopacity=.56\par
        ]
      \end{cookiecard}
  \end{columns}
\end{frame}
}

\begin{frame}{Show one layer at a time}{The geometry stays fixed across overlays}
  \begin{wide}
    \centering
    \begin{tikzpicture}[
      node distance=7mm and 10mm,
      endpoint/.style={cookie node,minimum width=1.75cm,minimum height=.78cm,font=\small\bfseries},
      system/.style={cookie node,minimum width=3.35cm,minimum height=1.02cm,font=\small\bfseries},
      rowlabel/.style={anchor=east,text=cookieMuted,font=\small\bfseries}
    ]
      % Audience-facing model
      \begin{scope}[muted on=<2-3>]
        \node[rowlabel] at (-5.05,1.72) {Your view};
        \node[endpoint] (shop1) at (-3.25,1.72) {Provider};
        \node[system] (magic) at (0,1.72) {\usebox{\cookieDemoPenguin}};
        \node[endpoint] (home1) at (3.25,1.72) {Home};
        \draw[cookie arrow] (shop1) -- (magic);
        \draw[cookie arrow] (magic) -- (home1);
      \end{scope}

      % Operator-facing model
      \begin{scope}[muted on=<{1,3}>]
        \node[rowlabel] at (-5.05,0) {System view};
        \node[endpoint] (shop2) at (-3.25,0) {Provider};
        \node[system] (dispatch) at (0,0) {Dispatch service};
        \node[endpoint] (home2) at (3.25,0) {Home};
        \draw[cookie arrow] ([yshift=2.2mm]shop2.east) -- ([yshift=2.2mm]dispatch.west);
        \draw[cookie arrow] (shop2.east) -- (dispatch.west);
        \draw[cookie arrow] ([yshift=-2.2mm]shop2.east) -- ([yshift=-2.2mm]dispatch.west);
        \draw[cookie arrow] ([yshift=2.2mm]dispatch.east) -- ([yshift=2.2mm]home2.west);
        \draw[cookie arrow] (dispatch.east) -- (home2.west);
        \draw[cookie arrow] ([yshift=-2.2mm]dispatch.east) -- ([yshift=-2.2mm]home2.west);
      \end{scope}

      % Worker-facing model
      \begin{scope}[muted on=<1-2>]
        \node[rowlabel] at (-5.05,-1.72) {Worker view};
        \node[endpoint] (shop3) at (-3.65,-1.72) {Provider};
        \node[endpoint] (queue) at (-1.05,-1.72) {Queue};
        \node[endpoint] (base) at (1.25,-1.72) {Base};
        \node[endpoint] (home3) at (3.85,-1.72) {Home};
        \draw[cookie arrow,draw=cookieLine,densely dashed] (shop3) -- (queue);
        \draw[cookie arrow] (queue) -- (base);
        \draw[cookie arrow,draw=cookieLine,densely dashed] (base) -- (home3);
      \end{scope}
    \end{tikzpicture}

    \vspace{.45em}
    \only<1>{\small Start with the audience's view of the system.}
    \only<2>{\small Then show the service boundary.}
    \only<3>{\small Add worker detail when it helps explain the failure mode.}
    \only<4>{\small End with all rows visible so the views can be compared.}
  \end{wide}

  \note{
    \begin{itemize}
      \item Advance once per row.
      \item Do not explain internal detail before the row appears.
      \item On overlay four, compare the rows.
      \item The penguin is Cookie's optional mascot and is available as \texttt{\string\cookiepenguin}.
    \end{itemize}
  }
\end{frame}

\subsection{Overlays}

\begin{frame}{Reveal without moving the layout}
  \begin{columns}[T,onlytextwidth]
    \column{.34\textwidth}
      \centering
      \resizebox{\linewidth}{!}{%
      \begin{tikzpicture}[
        node distance=7mm and 9mm,
        flow node/.style={
          cookie node,
          minimum width=1.74cm,
          minimum height=.72cm,
          font=\small\bfseries,
          align=center
        },
        flow active/.style={
          cookie accent node,
          minimum width=1.74cm,
          minimum height=.72cm,
          font=\small\bfseries,
          align=center
        },
        flow muted/.style={
          flow node,
          fill=cookieSoft!72!cookiePaper,
          text=cookieMuted
        },
        branch label/.style={
          font=\scriptsize\sffamily\bfseries,
          text=cookieMuted,
          inner sep=1pt
        }
      ]
        \node[flow active] (request) {Request};
        \node[flow active,below left=of request,visible on=<2->] (validate) {Validate};
        \node[flow muted,below right=of request,visible on=<2->] (fallback) {Fallback};
        \draw[cookie arrow,draw=cookieAccent,visible on=<2->]
          (request) -- node[branch label,above left] {primary} (validate);
        \draw[cookie arrow,draw=cookieMuted,visible on=<2->]
          (request) -- node[branch label,above right] {backup} (fallback);
        \node[flow active,below=10mm of validate,visible on=<3->] (readya) {Ready};
        \node[flow active,below=10mm of fallback,visible on=<3->] (readyb) {Ready};
        \draw[cookie arrow,draw=cookieAccent,visible on=<3->]
          (validate) -- node[branch label,left] {ok} (readya);
        \draw[cookie arrow,draw=cookieAccent,visible on=<3->]
          (fallback) -- node[branch label,right] {recover} (readyb);
        \draw[cookieAccent,line width=.8pt,densely dashed,visible on=<3->]
          (readya.east) -- (readyb.west);
        \node[
          below=5mm of $(readya)!0.5!(readyb)$,
          text=cookieAccent,
          font=\scriptsize\sffamily\bfseries,
          visible on=<3->
        ] {same final state};
      \end{tikzpicture}%
      }

    \column{.60\textwidth}
      \begin{itemize}
        \item<1-> \textcolor{cookieAccent}{Reserve the canvas.} The request node anchors every overlay.
        \item<2-> \textcolor{cookieAccent}{Reveal branches.} The primary path carries the accent; fallback stays muted.
        \item<3-> \textcolor{cookieAccent}{Show convergence.} Both paths end in the same user-visible state.
      \end{itemize}

      \vspace{.35em}
      \uncover<3->{%
        {\color{cookieLine}\rule{\linewidth}{.6pt}}\par\vspace{.7em}
        \begin{tabularx}{\linewidth}{@{}>{\bfseries}lX@{}}
          Primary & request $\rightarrow$ validate $\rightarrow$ ready\\
          Fallback & request $\rightarrow$ fallback $\rightarrow$ ready
        \end{tabularx}
      }
  \end{columns}

  \begin{modal}<4>[Overlay rule]
    \raggedright
    Hidden nodes are present from the first overlay. Each step changes opacity,
    so the accent path can appear without moving the text.
  \end{modal}
  \note{Advance through the request, branch, and convergence overlays. The modal explains why the layout does not jump.}
\end{frame}

\section{Content patterns}

\begin{frame}{Let the footer carry progress}
  \begin{wide}
    \centering
    \begin{tikzpicture}[node distance=1.55cm]
      \node[cookie accent node,minimum width=3.05cm,minimum height=1.2cm] (idea) {Idea};
      \node[cookie node,right=of idea,minimum width=3.05cm,minimum height=1.2cm] (evidence) {Evidence};
      \node[cookie node,right=of evidence,minimum width=3.05cm,minimum height=1.2cm] (decision) {Decision};
      \draw[cookie arrow] (idea) -- (evidence);
      \draw[cookie arrow] (evidence) -- (decision);
    \end{tikzpicture}

    \vspace{1.1em}
    The footer shows location and frame count; the slide can stay on the current step.
  \end{wide}
\end{frame}

\begin{frame}{Keep tables plain}
  \centering
  \begin{tabularx}{.82\textwidth}{@{}Xccc@{}}
    \toprule
    Feature & None & Simple & Progress \\
    \midrule
    Section page    & \cookiecheck & \cookiecheck & \cookiecheck \\
    Subsection page & \cookiecheck & \cookiecheck & \cookiecheck \\
    Progress bar    & head & title & foot \\
    Frame numbering & none & count & fraction \\
    \bottomrule
  \end{tabularx}

  \vspace{1em}
  \small Use \texttt{booktabs} for rules and \texttt{tabularx} for width. Cookie only sets spacing and color.
\end{frame}

\begin{frame}{A larger table}{Booktabs rules and aligned numeric columns}
  \centering
  \small
  \begin{tabularx}{\textwidth}{@{}Xrrr@{}}
    \toprule
    Faculty & With \TeX & Total & \% \\
    \midrule
    Informatics       & 1\,716 & 2\,904  & 59.09 \\
    Science           & 786    & 5\,275  & 14.90 \\
    Economics         & 64     & 4\,591  & 1.39  \\
    Arts              & 69     & 10\,000 & 0.69  \\
    Medicine          & 8      & 2\,014  & 0.40  \\
    Law               & 15     & 4\,824  & 0.31  \\
    Education         & 19     & 8\,219  & 0.23  \\
    Social studies    & 12     & 5\,599  & 0.21  \\
    Sports studies    & 3      & 2\,062  & 0.15  \\
    \bottomrule
  \end{tabularx}

  \vspace{.65em}
  \footnotesize Example data adapted from a university thesis-count table.
\end{frame}

\begin{frame}{A probability tree}{TikZ diagram adapted from the source example}
  \centering
  \tikzset{
    coin node/.style={
      draw=cookieInk,
      text=cookieInk,
      circle,
      minimum size=11mm,
      line width=.7pt,
      font=\Large\bfseries,
      inner sep=0pt
    },
    coin prob/.style={
      text=cookieInk,
      font=\Large\bfseries,
      inner sep=1pt
    },
    coin leafprob/.style={
      text=cookieInk,
      font=\Large\bfseries,
      inner sep=1pt
    }
  }
  \scalebox{.74}{%
  \begin{tikzpicture}[line cap=round,line join=round]
    \node[
      draw=cookieInk,
      text=cookieInk,
      minimum width=6.65cm,
      minimum height=.90cm,
      line width=.7pt,
      font=\Large\bfseries
    ] (title) at (0,2.42) {Coin flipping};

    \node[
      draw=cookieInk,
      text=cookieInk,
      shape=circle split,
      minimum size=13mm,
      line width=.7pt,
      font=\Large\bfseries,
      inner sep=0pt,
      rotate=30
    ] (start) at (0,1.10)
      {\rotatebox{-30}{H}\nodepart{lower}\rotatebox{-30}{T}};

    \node[coin node] (left) at (-2.35,-.22) {H};
    \node[coin node] (right) at (2.35,-.22) {T};
    \node[coin node] (hh) at (-3.45,-1.72) {H};
    \node[coin node] (ht) at (-1.25,-1.72) {T};
    \node[coin node] (th) at (1.25,-1.72) {H};
    \node[coin node] (tt) at (3.45,-1.72) {T};

    \draw[cookieInk,line width=.7pt] (title.south) -- (start.north);
    \draw[cookieInk,line width=.7pt] (start) -- (left);
    \draw[cookieInk,line width=.7pt] (start) -- (right);
    \draw[cookieInk,line width=.7pt] (left) -- (hh);
    \draw[cookieInk,line width=.7pt] (left) -- (ht);
    \draw[cookieInk,line width=.7pt] (right) -- (th);
    \draw[cookieInk,line width=.7pt] (right) -- (tt);

    \path (start) -- (left) node[coin prob,pos=.48,above left=1pt] {$0.5$};
    \path (start) -- (right) node[coin prob,pos=.48,above right=1pt] {$0.5$};
    \path (left) -- (hh) node[coin prob,pos=.42,above left=1pt] {$0.5$};
    \path (left) -- (ht) node[coin prob,pos=.42,above right=1pt] {$0.5$};
    \path (right) -- (th) node[coin prob,pos=.42,above left=1pt] {$0.5$};
    \path (right) -- (tt) node[coin prob,pos=.42,above right=1pt] {$0.5$};

    \node[coin leafprob] at (-3.45,-2.70) {$0.25$};
    \node[coin leafprob] at (-1.25,-2.70) {$0.25$};
    \node[coin leafprob] at (1.25,-2.70) {$0.25$};
    \node[coin leafprob] at (3.45,-2.70) {$0.25$};
  \end{tikzpicture}
  }
\end{frame}

\begin{frame}[standout,noframenumbering]
  One idea per slide.
\end{frame}

\begin{frame}{Use quotes sparingly}
  \begin{cookiequote}[Cookie design principle]
    A good slide should make the argument easier to see.
  \end{cookiequote}

  \vspace{.4em}
  A quote works best as a turn in the talk. Don't use too many!
\end{frame}

\begin{frame}[fragile]{Speaker notes live in the source}{Use \texttt{notes=second} for presenter mode}
  \begin{columns}[T,onlytextwidth]
    \column{.57\textwidth}
      \begin{cookiecode}[title=Inside any frame]{TeX}
The audience sees this.

\note{
  Remember why this slide matters.
  Add timing, transitions, and questions here.
}
      \end{cookiecode}

    \column{.39\textwidth}
      \begin{block}{Output modes}
        \small
        \begin{tabularx}{\linewidth}{@{}>{\ttfamily}lX@{}}
          notes=hide & slides only \\
          notes=show & slides plus notes \\
          notes=second & presenter layout \\
          notes=only & notes pages only
        \end{tabularx}
      \end{block}

      \cookiebadge{native Beamer} \hspace{.25em}\cookiebadge{styled notes page}
  \end{columns}

  \note{
    This note compiles even when the audience deck hides notes. Build with
    \texttt{notes=second} to see the presenter page.
  }
\end{frame}

\section{Configuration}

\begin{frame}{Options in one place}
  \begin{columns}[T,onlytextwidth]
    \column{.48\textwidth}
      \begin{block}{Layout}
        \small
        \begin{tabularx}{\linewidth}{@{}>{\ttfamily}lX@{}}
          progressbar & none, head, frametitle, foot \\
          numbering & none, counter, fraction \\
          numberrail & none, section, subsection \\
          covered & invisible, transparent \\
          notes & hide, show, second, only \\
          toc & none, aftertitle
        \end{tabularx}
      \end{block}

    \column{.48\textwidth}
      \begin{block}{Appearance}
        \small
        \begin{tabularx}{\linewidth}{@{}>{\ttfamily}lX@{}}
          accent & any xcolor name \\
          background & light, dark \\
          block & transparent, fill \\
          numberrailcolor & muted, accent \\
          titleformat & regular, smallcaps, allcaps
        \end{tabularx}
      \end{block}
  \end{columns}

  \vfill
  \small Compatibility aliases such as \texttt{secslide}, \texttt{subsecslide},
  \texttt{notoc}, and \texttt{coloraccent} make older awesome-beamer decks easier to move.
\end{frame}

\begin{frame}{Two background commands}
  \begin{columns}[T,onlytextwidth]
    \column{.48\textwidth}
      \begin{cookiecard}[title=Title-page wedge]
        \texttt{\textbackslash cookietitleimage\{image.jpg\}}

        \vspace{.45em}
        Clips the image into Cookie's asymmetric title art.
      \end{cookiecard}

    \column{.48\textwidth}
      \begin{cookiecard}[title=Ordinary frames]
        \texttt{bgimage=image.jpg}

        \vspace{.45em}
        Add the frame keys you need:
        \par\vspace{.25em}
        {\ttfamily\small
          bgfit=cover\par
          bgoverlay=cookieInk\par
          bgoverlayopacity=.56\par
        }
      \end{cookiecard}
  \end{columns}

  \vfill
  \begin{block}{Global background}
    \texttt{\textbackslash cookiebackgroundimage\{image.jpg\}} applies an image to later frames.\\
    \texttt{\textbackslash cookieclearbackgroundimage} removes it.
  \end{block}
\end{frame}

\begin{frame}{Build the demo with LuaLaTeX}
  \cookiekicker{Local build}
  \vspace{.55em}

  \begin{block}{Repository command}
    \texttt{latexmk demo.tex}
  \end{block}

  \begin{itemize}
    \item The local \texttt{.latexmkrc} selects LuaLaTeX and Biber for this demo.
    \item \texttt{latexmk} reruns TeX when section, frame, QR, or bibliography data changes.
    \item If you copy only the theme file into another deck, build that deck with LuaLaTeX.
  \end{itemize}
\end{frame}

\section{References}

\begin{frame}{Citations use normal BibLaTeX}{No theme-specific bibliography syntax}
  Cookie is still ordinary Beamer: the typesetting stack builds on \TeX\ \parencite{knuth1984},
  \LaTeX\ \parencite{lamport1994}, and Beamer \parencite{tantau2025}.

  \vspace{.8em}
  The design borrows progress bars from Metropolis \parencite{metropolis} and the numbered
  slide grammar of awesome-beamer \parencite{awesomebeamer}.

  \vfill
  \begin{block}{Source setup}
    \texttt{\textbackslash usepackage[backend=biber,style=authoryear]\{biblatex\}}\\
    \texttt{\textbackslash addbibresource\{refs.bib\}}
  \end{block}
\end{frame}

\begin{frame}[allowframebreaks]{References}
  \printbibliography[heading=none]
\end{frame}

\end{document}
